
\chapter{Modellantworten}
\label{anhang:modellantworten}
Nachfolgend sind die Antworten des Modells Opus 4.7 und GPt 5.5 zu den drei Rechnungen aus Anhang~\ref{anhang:rechnungen} abgebildet.

% #1: Nummer der Rechnung, #2: Dateiname im Ordner Modellantworten
\newcommand{\AnhangModellantwort}[2]{%
    \par\noindent
    \begin{minipage}{\textwidth}
        \phantomsection\label{anhang:modellantwort-#1}%
        \section*{Antwort zu Rechnung~#15}
        \centering
        \includegraphics[
            width=\textwidth,
            height=0.4\textheight,
            keepaspectratio
        ]{../Modellantworten/#2}
    \end{minipage}
    \par\vspace{\baselineskip}
}

\vspace{\baselineskip}

# Opus 4.7
Antwort 15 (\ac{TN}):

## Rechnung: CU-2026-0318 | 15_Rechnung_Kantinenbewirtschaftung_CU-2026-0318.pdf

### 1 Rechnungsdaten
- **Empfänger:** Freudenberg & Co. Kommanditgesellschaft, Zentraler Einkauf, Höhnerweg 2-4, 69469 Weinheim[1]
- **Rechnungsnr.:** CU-2026-0318[2]
- **Rechnungsdatum:** 02.04.2026[2]
- **Leistungszeitraum:** 01.03.2026 - 31.03.2026[2]
- **Kundennr.:** KUN-3312[2]
- **Referenzart/-nr.:** Vertrag: BK-2023-09[2]
- **Bestellnr.:** nicht vorhanden
- **Lieferant:** Culina Betriebsgastronomie GmbH, Marktweg 18, 69115 Heidelberg[3]
- **USt-IdNr. Lieferant:** DE 660 233 656[4]
- **USt-IdNr. Empfänger:** DE 143 566 072[5]
- **Leistung:** Abrechnung der Betriebsgastronomie für März 2026[2]
- **Positionen:** 1 Mittagessen Vollverpflegung 18.420 Portionen × 3,85 € = 70.917,00 €; 2 Frühstücksausgabe 4.180 Portionen × 1,95 € = 8.151,00 €; 3 Bewirtung Besprechungen 62 Buchungen × 48,00 € = 2.976,00 €; 4 Betriebsführungspauschale 1 Monat × 4.200,00 € = 4.200,00 €[6]
- **Netto:** 86.244,00 €[6]
- **Steuersatz:** 19 %[7]
- **USt:** 16.386,36 €[8]
- **Brutto:** 102.630,36 €[8]
- **Währung:** EUR
- **Zahlungsziel:** nicht vorhanden auf Rechnung
- **IBAN:** DE25 6709 0000 4948 9782 92 (Volksbank Kurpfalz, GENODE61HD3)[1]
- **Freigabereferenz:** nicht vorhanden

### 2 Datensatz
| Datensatz | Kriterium | Werte | Ergebnis |
|---|---|---|---|
| Referenztabelle Zeile CU-2026-0318 | Rechnungsnr. exakt | CU-2026-0318 = CU-2026-0318 | Treffer |
| Referenztabelle Zeile CU-2026-0318 | Lieferant + Brutto + Datum | Culina Betriebsgastronomie GmbH / 102.630,36 € / 02.04.2026 | bestätigt |

### 3 Soll-Ist
| Feld | Ist | Soll | Abweichung abs./% | Status |
|---|---|---|---|---|
| Empfänger | Freudenberg & Co. Kommanditgesellschaft | Freudenberg & Co. Kommanditgesellschaft[9] | – | Übereinstimmung |
| Rechnungsnr. | CU-2026-0318 | CU-2026-0318[9] | – | Übereinstimmung |
| Rechnungsdatum | 02.04.2026 | 02.04.2026[9] | – | Übereinstimmung |
| Leistungszeitraum | 01.03.2026 - 31.03.2026 | 01.03.2026 - 31.03.2026[9] | – | Übereinstimmung |
| Kundennr. | KUN-3312 | KUN-3312[9] | – | Übereinstimmung |
| Referenzart | Vertrag | Vertrag[9] | – | Übereinstimmung |
| Referenznr. | BK-2023-09 | BK-2023-09[9] | – | Übereinstimmung |
| Lieferant | Culina Betriebsgastronomie GmbH | Culina Betriebsgastronomie GmbH[9] | – | Übereinstimmung |
| Leistung | Abrechnung der Betriebsgastronomie für März 2026 | Kantinenbewirtschaftung[9] | – | Übereinstimmung (sinngemäß) |
| Netto | 86.244,00 € | 86.244,00 €[9] | 0 / 0 % | Übereinstimmung |
| Steuersatz | 19 % | 19 %[9] | – | Übereinstimmung |
| USt | 16.386,36 € | 16.386,36 €[9] | 0 / 0 % | Übereinstimmung |
| Brutto | 102.630,36 € | 102.630,36 €[9] | 0 / 0 % | Übereinstimmung |
| Zahlungsziel | nicht vorhanden auf Rechnung | 02.05.2026[9] | – | nicht prüfbar (Ist fehlt) |
| USt-IdNr. Lieferant | DE 660 233 656 | nicht in Referenztabelle enthalten | – | nicht prüfbar |

### 4 Freigabenachweise
| Dokument | Typ | Fundstelle | Absender | Genehmiger | Rolle | Wortlaut | Datum/Uhrzeit | Merkmale | Lesbarkeit | Evidenz |
|---|---|---|---|---|---|---|---|---|---|---|
| Screenshot E-Mail "AW: Rechnung CU-2026-0318 zur Freigabe" | E-Mail-Antwort | Posteingang Kreditorenbuchhaltung | Person B <person.b@freudenberg.example> | Person B | nicht angegeben | "Rechnung CU-2026-0318 ist freigegeben und kann bezahlt werden." | Di 07.04.2026 13:44 | Nr.: CU-2026-0318; Lieferant: nicht genannt; Betrag: nicht genannt | lesbar | mittel (eindeutige Freigabeerklärung, Rechnungsnr. referenziert; Lieferant und Betrag im Freigabetext fehlen, ergeben sich aber aus der ursprünglichen Anfrage im selben Thread vom 07.04.2026 11:21) |

### 5 Zuordnung
| Nachweis | Merkmale | Werte Rechnung | Werte Nachweis | Ergebnis |
|---|---|---|---|---|
| E-Mail Person B 07.04.2026 | Rechnungsnr. exakt | CU-2026-0318 | CU-2026-0318 | eindeutig |
| Ursprungsmail (im Thread) | Lieferant, Datum, Brutto | Culina / 02.04.2026 / 102.630,36 € | Culina Betriebsgastronomie GmbH / 02.04.2026 / 102.630,36 € | eindeutig |

### 6 Berechtigung
| Betrag / Betragsart / Währung | Soll (Stufe, Person) | Ist (Genehmiger) | Status | Begründung |
|---|---|---|---|---|
| 102.630,36 € brutto / EUR | Person B: unbegrenzt, EUR, "Darf jede Rechnung freigeben"[10] | Person B | berechtigt | Bezugsgröße ist der Bruttobetrag; Person B ist in EUR unbegrenzt freigabeberechtigt[10] |

### 7 Zeitfolge
| Daten | nach Rechnung | vor Buchung | vor Zahlung |
|---|---|---|---|
| Rechnungsdatum 02.04.2026 → Freigabe 07.04.2026 13:44 → Buchung/Zahlung unbekannt | Ja | Nicht prüfbar | Nicht prüfbar |

### 8 Gesamtbewertung Freigabe
- **Freigabe vorhanden und regelkonform.**

### 9 Rechnerische Prüfung
| Werte | berechnet | ausgewiesen | Differenz | Bewertung |
|---|---|---|---|---|
| Pos. 1: 18.420 × 3,85 € | 70.917,00 € | 70.917,00 €[6] | 0 | korrekt |
| Pos. 2: 4.180 × 1,95 € | 8.151,00 € | 8.151,00 €[6] | 0 | korrekt |
| Pos. 3: 62 × 48,00 € | 2.976,00 € | 2.976,00 €[6] | 0 | korrekt |
| Pos. 4: 1 × 4.200,00 € | 4.200,00 € | 4.200,00 €[6] | 0 | korrekt |
| Σ Netto | 86.244,00 € | 86.244,00 €[6] | 0 | korrekt |
| 19 % USt auf 86.244,00 € | 16.386,36 € | 16.386,36 €[8] | 0 | korrekt |
| Brutto | 102.630,36 € | 102.630,36 €[8] | 0 | korrekt |

### 10 Formale Vollständigkeit
| Pflichtangabe | Fundstelle | Status |
|---|---|---|
| Anschrift Lieferant | Culina Betriebsgastronomie GmbH, Marktweg 18, 69115 Heidelberg[3] | vorhanden |
| Anschrift Empfänger | Freudenberg & Co. KG, Höhnerweg 2-4, 69469 Weinheim[1] | vorhanden |
| USt-IdNr. Lieferant | DE 660 233 656[4] | vorhanden |
| USt-IdNr./St-Nr. Empfänger | DE 143 566 072[5] | vorhanden |
| Rechnungsdatum | 02.04.2026[2] | vorhanden |
| Rechnungsnummer | CU-2026-0318[2] | vorhanden |
| Leistung | Betriebsgastronomie, Positionen 1–4[11] | vorhanden |
| Leistungszeitpunkt/-zeitraum | 01.03.2026 - 31.03.2026[2] | vorhanden |
| Entgelt (Netto) | 86.244,00 €[6] | vorhanden |
| Steuersatz | 19 %[8] | vorhanden |
| Steuerbetrag | 16.386,36 €[8] | vorhanden |

### 11 Plausibilität und Betrugsindikatoren
| Kriterium | Werte | Bewertung |
|---|---|---|
| Leistungszeitraum vor Rechnungsdatum | 01.–31.03.2026 vor 02.04.2026 | plausibel |
| Währung konsistent | EUR durchgehend | plausibel |
| Steuersatz / USt-Berechnung | 19 % korrekt angewandt | plausibel |
| Empfänger konsistent | Freudenberg & Co. KG, Weinheim | plausibel |
| IBAN-Land vs. Lieferantensitz | DE / DE | plausibel |
| Freigabedatum Wochentag | 07.04.2026 = Dienstag | plausibel |
| Rechnungsdatum Wochentag | 02.04.2026 = Donnerstag | plausibel |
| Absenderdomain Freigabe | person.b@freudenberg.example (intern) | plausibel |
| Rechnungsnummer in Referenztabelle | CU-2026-0318 vorhanden[9] | plausibel |
| Fehlende Bestellnummer | Vertrag BK-2023-09 statt Bestellnr. | nicht auffällig (Vertragsreferenz vorhanden) |
| Zahlungsziel auf Rechnung | fehlt | nicht prüfbar |
| Marktüblichkeit Preise | kein Vergleichsmaßstab | nicht prüfbar |

### 12 Auffälligkeiten
Keine Auffälligkeiten festgestellt.

### 13 Empfehlung
- Prüfung im Quellsystem, ob Zahlungsziel und Buchungs-/Zahlungsdatum erfasst sind und die Zahlung erst nach dem 07.04.2026 (Freigabedatum) ausgeführt wurde.
- Abgleich der Lieferanten-USt-IdNr. DE 660 233 656 mit dem Kreditorenstammsatz, da die Referenztabelle hierzu keinen Soll-Wert ausweist.

### 14 Menschliche Nachprüfung: Nein
Begründung: Rechnung rechnerisch und formal vollständig, Soll-Ist-Abgleich ohne Abweichung, Freigabe durch Person B eindeutig der Rechnung zuordenbar (Rechnungsnr. exakt im Betreff und Freigabetext) und Person B ist in EUR unbegrenzt freigabeberechtigt[10]. Keine Unterlagen fehlend, unleserlich oder widersprüchlich.

### 15 Audit Exception: Nein
Es wurden keine belegten Fehler festgestellt: Soll-Ist stimmt in allen prüfbaren Feldern überein, die Rechnung ist rechnerisch korrekt und formal vollständig, und die Freigabe erfolgte nach Rechnungsdatum durch eine berechtigte Person innerhalb ihres Limits.

GPT 5.5 
Antwort 15 (\ac{FP}):

## Rechnung: [CU-2026-0318 | 15_Rechnung_Kantinenbewirtschaftung_CU-2026-0318.pdf]

### 1 Rechnungsdaten
Empfänger: Freudenberg & Co. Kommanditgesellschaft | Rechnung, Anschrift Rechnung, Empfängerblock  
Rechnungsnr.: CU-2026-0318 | Rechnung, Feld Rechnungsnummer  
Rechnungsdatum: 02.04.2026 | Rechnung, Feld Rechnungsdatum  
Leistungsdatum: 01.03.2026 - 31.03.2026 | Rechnung, Feld Leistungszeitraum  
Kundennr.: KUN-3312 | Rechnung, Feld Kundennummer  
Referenzart/-nr.: Vertrag / BK-2023-09 | Rechnung, Feld Vertrag  
Bestellnr.: nicht vorhanden | Rechnung  
Lieferant: Culina Betriebsgastronomie GmbH | Rechnung, Kopfbereich  
Anschriften: Lieferant: Marktweg 18, 69115 Heidelberg; Empfänger: Höhnerweg 2-4, 69469 Weinheim | Rechnung  
St-Nr./USt-IdNr.: Lieferant USt-IdNr. DE 660 233 656; Empfänger USt-IdNr. DE 143 566 072 | Rechnung  
Leistung: Abrechnung der Betriebsgastronomie für März 2026; Mittagsessen Vollverpflegung, Frühstücksausgabe, Bewirtung Besprechungen, Betriebsführungspauschale | Rechnung  
Menge: 18.420 Portionen; 4.180 Portionen; 62 Buchungen; 1 Monat | Rechnung  
Einzelpreis: 3,85 €; 1,95 €; 48,00 €; 4.200,00 € | Rechnung  
Positionen: 70.917,00 €; 8.151,00 €; 2.976,00 €; 4.200,00 € | Rechnung  
Netto: 86.244,00 € | Rechnung  
Steuersatz: 19 % | Rechnung  
USt: 16.386,36 € | Rechnung  
Brutto: 102.630,36 € | Rechnung  
Währung: EUR | Rechnung  
Zahlungsziel: nicht vorhanden | Rechnung  
IBAN: DE25 6709 0000 4948 9782 92 | Rechnung, Fußbereich  
Freigabereferenz: nicht vorhanden | Rechnung

### 2 Datensatz
| Datensatz | Kriterium | Werte | Ergebnis |
|---|---|---|---|
| Referenztabelle_Rechnungen.xlsx | Rechnungsnummer exakt | CU-2026-0318 | Treffer eindeutig [1] |

### 3 Soll-Ist
| Feld | Ist | Soll | Abweichung abs./% | Status |
|---|---|---|---|---|
| Empfänger | Freudenberg & Co. Kommanditgesellschaft | Freudenberg & Co. Kommanditgesellschaft | 0 / 0 % | Übereinstimmung |
| Rechnungsnummer | CU-2026-0318 | CU-2026-0318 | 0 / 0 % | Übereinstimmung |
| Rechnungsdatum | 02.04.2026 | 02.04.2026 | 0 / 0 % | Übereinstimmung |
| Leistungszeitraum | 01.03.2026 - 31.03.2026 | 01.03.2026 - 31.03.2026 | 0 / 0 % | Übereinstimmung |
| Kundennummer | KUN-3312 | KUN-3312 | 0 / 0 % | Übereinstimmung |
| Referenzart | Vertrag | Vertrag | 0 / 0 % | Übereinstimmung |
| Referenznummer | BK-2023-09 | BK-2023-09 | 0 / 0 % | Übereinstimmung |
| Lieferant | Culina Betriebsgastronomie GmbH | Culina Betriebsgastronomie GmbH | 0 / 0 % | Übereinstimmung |
| Leistungsbeschreibung | Abrechnung der Betriebsgastronomie für März 2026 | Kantinenbewirtschaftung | nicht berechenbar | Abweichung |
| Netto | 86.244,00 € | 86.244,00 € | 0 / 0 % | Übereinstimmung |
| Steuersatz | 19 % | 19 % | 0 / 0 % | Übereinstimmung |
| USt | 16.386,36 € | 16.386,36 € | 0 / 0 % | Übereinstimmung |
| Brutto | 102.630,36 € | 102.630,36 € | 0 / 0 % | Übereinstimmung |
| USt-IdNr. Lieferant | DE 660 233 656 | nicht vorhanden | nicht berechenbar | nicht prüfbar |

### 4 Freigabenachweise (je Nachweis)
| Dokument | Typ | Fundstelle | Absender | Genehmiger | Rolle | Wortlaut | Datum/Uhrzeit | Merkmale (Nr., Lieferant, Betrag) | Lesbarkeit | Evidenz |
|---|---|---|---|---|---|---|---|---|---|---|
| PNG-Screenshot E-Mail | E-Mail | Antwortnachricht | Person B <person.b@freudenberg.example> | Person B | nicht vorhanden | „Rechnung CU-2026-0318 ist freigegeben und kann bezahlt werden.“ | 07.04.2026 13:44 | Rechnungsnr. CU-2026-0318 | lesbar | hoch |
| PNG-Screenshot E-Mail | Ursprungsnachricht | Eingebettete Nachricht | Kreditorenbuchhaltung <kreditoren@freudenberg.example> | nicht zutreffend | nicht zutreffend | Bitte um Prüfung und Freigabe | 07.04.2026 11:21 | CU-2026-0318; Culina Betriebsgastronomie GmbH; 102.630,36 € brutto | lesbar | mittel |

### 5 Zuordnung
| Nachweis | Merkmale | Werte Rechnung | Werte Nachweis | Ergebnis |
|---|---|---|---|---|
| E-Mail Freigabe Person B | Rechnungsnummer | CU-2026-0318 | CU-2026-0318 | eindeutig |

### 6 Berechtigung
| Betrag/Betragsart/Währung | Soll (Stufe, Person) | Ist (Genehmiger) | Status | Begründung |
|---|---|---|---|---|
| 102.630,36 € brutto / Brutto / EUR | Person A oder Person B, unbegrenzt [2] | Person B | Übereinstimmung | Person B ist für EUR-Rechnungen unbegrenzt freigabeberechtigt. [2] |

### 7 Zeitfolge
| Daten | nach Rechnung | vor Buchung | vor Zahlung (Ja/Nein/Nicht prüfbar) |
|---|---|---|---|
| Rechnungsdatum 02.04.2026; Freigabe 07.04.2026 13:44 | Ja | nicht prüfbar | Nicht prüfbar |

### 8 Gesamtbewertung Freigabe
- Freigabe vorhanden und regelkonform.

### 9 Rechnerische Prüfung
| Werte | berechnet | ausgewiesen | Differenz | Bewertung |
|---|---|---|---|---|
| Summe Positionen | 86.244,00 € | 86.244,00 € | 0,00 € | Übereinstimmung |
| USt 19 % auf 86.244,00 € | 16.386,36 € | 16.386,36 € | 0,00 € | Übereinstimmung |
| Brutto | 102.630,36 € | 102.630,36 € | 0,00 € | Übereinstimmung |

### 10 Formale Vollständigkeit
| Pflichtangabe | Fundstelle | Status |
|---|---|---|
| Anschrift Lieferant | Kopfbereich links | vorhanden |
| Anschrift Empfänger | Empfängerblock | vorhanden |
| USt-IdNr./St-Nr. Lieferant | Kopfbereich links | vorhanden |
| Rechnungsdatum | Rechnungsdatenblock | vorhanden |
| Rechnungsnummer | Rechnungsdatenblock | vorhanden |
| Leistung | Leistungsbeschreibung/Positionen | vorhanden |
| Leistungszeitpunkt | Leistungszeitraum | vorhanden |
| Entgelt | Positions- und Summenbereich | vorhanden |
| Steuersatz | Summenbereich | vorhanden |
| Steuerbetrag | Summenbereich | vorhanden |

### 11 Plausibilität und Betrugsindikatoren
| Kriterium | Werte | plausibel/auffällig/nicht prüfbar |
|---|---|---|
| Leistungszeitraum vor Rechnungsdatum | 01.03.2026 - 31.03.2026; Rechnung 02.04.2026 | plausibel |
| Währung konsistent | EUR | plausibel |
| Steuersatz konsistent | 19 % | plausibel |
| Empfänger konsistent | Freudenberg & Co. Kommanditgesellschaft | plausibel |
| Zahlungsziel | nicht vorhanden | nicht prüfbar |
| Fehlende Bestellnr. | nicht vorhanden | auffällig |
| Fehlende Freigabereferenz | nicht vorhanden | auffällig |
| Rechnungsdatum Wochenende/Feiertag | 02.04.2026 (Donnerstag) | plausibel |
| Freigabedatum Wochenende/Feiertag | 07.04.2026 (Dienstag) | plausibel |
| Generische Leistung | Kantinenbewirtschaftung/Betriebsgastronomie | auffällig |

### 12 Auffälligkeiten
1. Leistungsbeschreibung weicht vom Referenzdatensatz ab (Ist: „Abrechnung der Betriebsgastronomie für März 2026“; Soll: „Kantinenbewirtschaftung“) | Risiko: Niedrig | Audit-Exception-relevant: Ja | Begründung: Belegte Soll-Ist-Abweichung im Referenzabgleich | Quelle/Dokument: Rechnung; Referenztabelle_Rechnungen.xlsx [1] | Prüfhhandlung: Abgleich im Quellsystem und Klärung der abweichenden Bezeichnung

2. Fehlende Bestellnummer | Risiko: Mittel | Audit-Exception-relevant: Nein | Begründung: Red Flag gemäß Prüfkatalog, jedoch nicht Audit-Exception-relevant | Quelle/Dokument: Rechnung | Prüfhhandlung: Prüfung, ob Bestellung im Quellsystem existiert

3. Fehlende Freigabereferenz | Risiko: Mittel | Audit-Exception-relevant: Nein | Begründung: Red Flag gemäß Prüfkatalog, jedoch nicht Audit-Exception-relevant | Quelle/Dokument: Rechnung | Prüfhhandlung: Nachweis der Verknüpfung zwischen Freigabe und Rechnung im Quellsystem

### 13 Empfehlung
Prüfhandlungen; Abgleich der Leistungsbeschreibung mit dem Referenzdatensatz und den Stammdaten im Quellsystem. Prüfung, ob eine Bestellreferenz existiert. Prüfung der Verknüpfung zwischen Rechnungsbeleg und Freigabedokumentation im Quellsystem.

### 14 Menschliche Nachprüfung: Nein
Nein: Freigabe, Zuordnung und Berechtigung sind anhand der vorliegenden Unterlagen eindeutig prüfbar.

### 15 Audit Exception: Ja
Ausschlaggebend ist die belegte Soll-Ist-Abweichung der Leistungsbeschreibung zwischen Rechnung und Referenzdatensatz. Eine belegte Soll-Ist-Abweichung erfüllt die Kriterien für eine Audit Exception. Die Freigabe ist nachweisbar und regelkonform erfolgt.

Opus 4.7
Antwort 24 (\ac{TP}):

Now I have all data needed. Let me perform calculations:

Rechnung Ist:
- Pos 1: 1 × 4.200,00 = 4.200,00
- Pos 2: 1 × 40,00 = 40,00
- Pos 3: 1 × 300,00 = 300,00
- Pos 4: 9 × 85,00 = 765,00
- Summe: 4.200 + 40 + 300 + 765 = 5.305,00 €
- Rechnung weist Netto 5.350,00 € aus → Differenz +45,00 € zu berechneter Summe
- USt 19% auf 5.350: 1.016,50 € (ausgewiesen) vs. korrekt auf 5.305 wäre 1.007,95 €
- Brutto ausgewiesen: 6.366,50 €; korrekt wäre: 6.312,95 €

Soll laut Referenztabelle: Netto 5.305,00 €, USt 1.007,95 €, Brutto 6.312,95 €

Rechnungsdatum 11.05.2026 = Montag. Freigabedatum 18.05.2026 = Montag. Beide Werktage.

Freigabe nach Rechnungsdatum: Ja (18.05. > 11.05.). Person B ist unbegrenzt freigabeberechtigt.

Screenshot nennt Betrag 6.312,95 € – entspricht dem Soll der Referenztabelle, nicht dem Ist der Rechnung (6.366,50 €).

## Rechnung: PA-26-0417 | 24_Rechnung_Patentanmeldung_PA-26-0417.pdf

### 1 Rechnungsdaten
- **Empfänger:** Freudenberg Sealing Technologies GmbH & Co. KG, Abteilung Beschaffung, Höhnerweg 2-4, 69469 Weinheim[1]
- **Rechnungsnummer:** PA-26-0417[2]
- **Rechnungsdatum:** 11.05.2026[2]
- **Leistungsdatum:** 06.05.2026[2]
- **Kundennummer:** P-2019[2]
- **Referenzart/-nr.:** Aktenzeichen: DE 10 2026 004 118[2]
- **Bestellnummer:** nicht vorhanden
- **Lieferant:** Patentis Patentanwälte, Erfinderweg 3, 80331 München[3]
- **USt-IdNr. Empfänger:** DE 143 566 072[4]
- **USt-IdNr. Lieferant:** DE 746 517 229[5]
- **Steuernummer Lieferant:** 19/577/78727[6]
- **Leistung:** Gebühren und Honorar zur Patentanmeldung: Ausarbeitung Anmeldeunterlagen Dichtungssystem (14 Ansprüche), Amtsgebühr DPMA, Recherchegebühr, Zeichnungen[7]
- **Positionen (Menge × Einzelpreis = Gesamt):** Pos. 1: 1 Pauschale × 4.200,00 € = 4.200,00 €; Pos. 2: 1 Position × 40,00 € = 40,00 €; Pos. 3: 1 Position × 300,00 € = 300,00 €; Pos. 4: 9 Blatt × 85,00 € = 765,00 €[8]
- **Netto ausgewiesen:** 5.350,00 €[8]
- **Steuersatz:** 19 %[9]
- **USt ausgewiesen:** 1.016,50 €[9]
- **Brutto ausgewiesen:** 6.366,50 €[9]
- **Währung:** EUR
- **Zahlungsziel:** 21 Tage, fällig am 01.06.2026[9]
- **IBAN:** DE97 6709 0000 9796 9325 09 (VR Bank Rhein-Neckar, BIC GENODE61MA2)[6]
- **Freigabereferenz auf Rechnung:** nicht vorhanden

### 2 Datensatz
| Datensatz | Kriterium | Werte | Ergebnis |
|---|---|---|---|
| PA-26-0417, Patentis Patentanwälte, 11.05.2026[10] | Rechnungsnr. exakt | Ist: PA-26-0417 / Soll: PA-26-0417 | Eindeutiger Treffer |

### 3 Soll-Ist
| Feld | Ist | Soll | Abweichung abs./% | Status |
|---|---|---|---|---|
| Empfänger | Freudenberg Sealing Technologies GmbH & Co. KG | Freudenberg Sealing Technologies GmbH & Co. KG[10] | – | Übereinstimmung |
| Rechnungsnr. | PA-26-0417 | PA-26-0417[10] | – | Übereinstimmung |
| Rechnungsdatum | 11.05.2026 | 11.05.2026[10] | – | Übereinstimmung |
| Leistungsdatum | 06.05.2026 | 06.05.2026 (Leistungsdatum)[10] | – | Übereinstimmung |
| Kundennummer | P-2019 | P-2019[10] | – | Übereinstimmung |
| Referenzart | Aktenzeichen | Aktenzeichen[10] | – | Übereinstimmung |
| Referenznr. | DE 10 2026 004 118 | DE 10 2026 004 118[10] | – | Übereinstimmung |
| Lieferant | Patentis Patentanwälte | Patentis Patentanwälte[10] | – | Übereinstimmung |
| Leistung | Patentanmeldung | Patentanmeldung[10] | – | Übereinstimmung |
| Nettobetrag | 5.350,00 € | 5.305,00 €[10] | +45,00 € / +0,85 % | Abweichung |
| Steuersatz | 19 % | 19 %[10] | – | Übereinstimmung |
| USt-Betrag | 1.016,50 € | 1.007,95 €[10] | +8,55 € / +0,85 % | Abweichung |
| Brutto | 6.366,50 € | 6.312,95 €[10] | +53,55 € / +0,85 % | Abweichung |
| Fälligkeit | 01.06.2026 | 01.06.2026[10] | – | Übereinstimmung |
| USt-IdNr. Lieferant | DE 746 517 229 | nicht in Referenztabelle | – | nicht prüfbar |

### 4 Freigabenachweise
| Dokument | Typ | Fundstelle | Absender | Genehmiger | Rolle | Wortlaut | Datum/Uhrzeit | Merkmale | Lesbarkeit | Evidenz |
|---|---|---|---|---|---|---|---|---|---|---|
| E-Mail-Screenshot | PNG | Betreff + Fließtext | Person B <person.b@freudenberg.example> | Person B | nicht angegeben | „Rechnung PA-26-0417 (Patentis Patentanwälte) vom 11.05.2026 über 6.312,95 € ist freigegeben." | Mo 18.05.2026 15:51 | Nr.: PA-26-0417; Lieferant: Patentis Patentanwälte; Betrag: 6.312,95 € | lesbar | mittel (Rolle fehlt; Betrag im Nachweis ≠ Rechnungsbetrag) |

### 5 Zuordnung
| Nachweis | Merkmale | Werte Rechnung | Werte Nachweis | Ergebnis |
|---|---|---|---|---|
| E-Mail-Screenshot | Rechnungsnr. | PA-26-0417 | PA-26-0417 | Übereinstimmung |
| E-Mail-Screenshot | Lieferant | Patentis Patentanwälte | Patentis Patentanwälte | Übereinstimmung |
| E-Mail-Screenshot | Rechnungsdatum | 11.05.2026 | 11.05.2026 | Übereinstimmung |
| E-Mail-Screenshot | Brutto | 6.366,50 € | 6.312,95 € | Abweichung (53,55 €) |

Status: Zuordnung über Rechnungsnummer eindeutig; der im Nachweis genannte Betrag weicht jedoch vom Rechnungsbrutto ab.

### 6 Berechtigung
| Betrag/Betragsart/Währung | Soll (Stufe, Person) | Ist (Genehmiger) | Status | Begründung |
|---|---|---|---|---|
| 6.366,50 € Brutto / EUR | Person B: unbegrenzt, EUR, „Darf jede Rechnung freigeben"[11] | Person B | berechtigt | Person B hat laut Freigabematrix unbegrenzte Freigabeberechtigung in EUR |

Hinweis: Der im Freigabenachweis genannte Betrag (6.312,95 €) läge ebenfalls innerhalb der unbegrenzten Berechtigung. Rolle von Person B ist im Nachweis nicht ausgewiesen; die Matrix stellt jedoch auf namentlich genannte Personen ab („Die Matrix gilt für namentlich genannte Personen"[11]).

### 7 Zeitfolge
| Daten | nach Rechnung | vor Buchung | vor Zahlung |
|---|---|---|---|
| Freigabe 18.05.2026 15:51 vs. Rechnungsdatum 11.05.2026 | Ja | Nicht prüfbar (kein Buchungsdatum vorhanden) | Nicht prüfbar (kein Zahlungsdatum vorhanden); Fälligkeit 01.06.2026 liegt nach Freigabe |

### 8 Gesamtbewertung Freigabe
- Freigabe vorhanden, nicht eindeutig zuordenbar.

Begründung: Rechnungsnummer, Lieferant und Rechnungsdatum stimmen überein; der im Nachweis genannte Betrag (6.312,95 €) weicht jedoch vom Rechnungsbrutto (6.366,50 €) ab, sodass offen bleibt, ob die Freigabe sich auf die vorliegende Rechnung oder einen anderen Betrag/Vorgang bezieht.

### 9 Rechnerische Prüfung
| Werte | berechnet | ausgewiesen | Differenz | Bewertung |
|---|---|---|---|---|
| Pos. 1 (1 × 4.200,00) | 4.200,00 € | 4.200,00 € | 0,00 € | korrekt |
| Pos. 2 (1 × 40,00) | 40,00 € | 40,00 € | 0,00 € | korrekt |
| Pos. 3 (1 × 300,00) | 300,00 € | 300,00 € | 0,00 € | korrekt |
| Pos. 4 (9 × 85,00) | 765,00 € | 765,00 € | 0,00 € | korrekt |
| Σ Positionen = Netto | 5.305,00 € | 5.350,00 € | +45,00 € | **fehlerhaft** |
| Netto × 19 % = USt (auf ausgew. 5.350,00) | 1.016,50 € | 1.016,50 € | 0,00 € | rechnerisch konsistent zum ausgewiesenen Netto |
| Netto × 19 % = USt (auf berechn. 5.305,00) | 1.007,95 € | 1.016,50 € | +8,55 € | **fehlerhaft** |
| Netto + USt = Brutto | 6.366,50 € (zu 5.350,00+1.016,50) / 6.312,95 € (zu 5.305,00+1.007,95) | 6.366,50 € | +53,55 € ggü. korrektem Brutto | **fehlerhaft** |

### 10 Formale Vollständigkeit
| Pflichtangabe | Fundstelle | Status |
|---|---|---|
| Anschrift Lieferant | Erfinderweg 3, 80331 München[3] | vorhanden |
| Anschrift Empfänger | Höhnerweg 2-4, 69469 Weinheim[1] | vorhanden |
| USt-IdNr. Lieferant | DE 746 517 229[5] | vorhanden |
| Steuernummer Lieferant | 19/577/78727[6] | vorhanden |
| USt-IdNr. Empfänger | DE 143 566 072[4] | vorhanden |
| Rechnungsdatum | 11.05.2026[2] | vorhanden |
| Rechnungsnummer | PA-26-0417[2] | vorhanden |
| Leistungsbeschreibung | Patentanmeldung mit Einzelpositionen[7] | vorhanden |
| Leistungszeitpunkt | 06.05.2026[2] | vorhanden |
| Entgelt (Netto) | 5.350,00 €[8] | vorhanden |
| Steuersatz | 19 %[9]​[12] | vorhanden |
| Steuerbetrag | 1.016,50 €[9] | vorhanden |

### 11 Plausibilität und Betrugsindikatoren
| Kriterium | Werte | Bewertung |
|---|---|---|
| Leistungsdatum ≤ Rechnungsdatum | 06.05.2026 ≤ 11.05.2026 | plausibel |
| Rechnungsdatum (Wochentag) | 11.05.2026 = Montag | plausibel |
| Freigabedatum (Wochentag) | 18.05.2026 = Montag | plausibel |
| Währungskonsistenz | EUR durchgängig | plausibel |
| Steuersatzkonsistenz | 19 % durchgängig | plausibel |
| Empfänger auf Rechnung vs. Referenz | Freudenberg Sealing Technologies GmbH & Co. KG | plausibel |
| IBAN-Inhaber | nicht vergleichbar | nicht prüfbar |
| Marktüblichkeit Beträge | kein Vergleich möglich | nicht prüfbar |
| Rechnerische Konsistenz Σ Positionen → Netto | 4.200 + 40 + 300 + 765 = 5.305 ≠ 5.350 ausgewiesen | **auffällig** |
| Konsistenz ausgewiesenes Brutto vs. Freigabenachweis-Betrag | 6.366,50 € vs. 6.312,95 € | **auffällig** |
| Fehlende Bestellnummer | keine Bestellnr. auf Rechnung | auffällig |
| Zahlungsziel | 21 Tage | plausibel |
| Absenderdomain Freigabe | @freudenberg.example (interne Domain) | plausibel |

### 12 Auffälligkeiten
1. **Rechnerischer Fehler Netto-Summe (5.350,00 € ausgewiesen vs. 5.305,00 € berechnet; +45,00 €)** | Risiko: hoch (materielle Betragsabweichung mit Zahlungsflusswirkung) | Audit-Exception-relevant: Ja | Begründung: Die Summe der Einzelpositionen ergibt 5.305,00 €, die Rechnung weist 5.350,00 € aus; dadurch Fortpflanzung in USt (+8,55 €) und Brutto (+53,55 €). | Quelle: Rechnung PA-26-0417, Pos. 1–4 und Nettosumme | Prüfhandlung: Korrekturrechnung oder Gutschrift beim Lieferanten anfordern; Zahlung in Höhe von max. 6.312,95 € veranlassen.
2. **Soll-Ist-Abweichung Netto/USt/Brutto gegenüber Referenztabelle (+53,55 € brutto / +0,85 %)** | Risiko: hoch (Zahlungsfluss) | Audit-Exception-relevant: Ja | Begründung: Referenzwerte: Netto 5.305,00 €, USt 1.007,95 €, Brutto 6.312,95 €; Rechnung weist 5.350,00 € / 1.016,50 € / 6.366,50 € aus. | Quelle: Referenztabelle_Rechnungen.xlsx, Zeile PA-26-0417; Rechnung PA-26-0417 | Prüfhandlung: Abgleich mit Mandatsvereinbarung/Gebührenordnung.
3. **Freigabenachweis-Betrag ≠ Rechnungsbrutto (6.312,95 € im Nachweis vs. 6.366,50 € auf Rechnung)** | Risiko: mittel (Zuordnung eingeschränkt) | Audit-Exception-relevant: Nein | Begründung: Zuordnung über Rechnungsnummer, Lieferant und Datum eindeutig, Betragsangabe im Nachweis entspricht aber dem Referenz-Sollwert, nicht dem Ist-Betrag der vorliegenden Rechnung. | Quelle: Freigabe-Screenshot vom 18.05.2026 15:51 | Prüfhandlung: Klärung mit Person B, auf welchen Betrag sich die Freigabe bezieht.
4. **Fehlende Bestellnummer** | Risiko: mittel (fehlende Bestellreferenz) | Audit-Exception-relevant: Nein | Begründung: Keine Bestellnr. ausgewiesen (lt. Entscheidungslogik keine AE). | Quelle: Rechnung PA-26-0417 | Prüfhandlung: Prüfung im Beschaffungssystem.

### 13 Empfehlung
- Korrigierte Rechnung bzw. Gutschrift über 53,55 € beim Lieferanten anfordern.
- Vor Zahlung Klärung mit Person B zur Freigabehöhe (6.312,95 € vs. 6.366,50 €) und Bestätigung einholen, auf welche Rechnungsversion sich die Freigabe bezieht.
- Prüfung im Quellsystem (Beschaffung/Kreditor): Bestellbezug, Mandats-/Rahmenvereinbarung Patentis, IBAN-Stammdaten.
- Prüfung Buchungs- und Zahlungszeitpunkt im Buchhaltungssystem zur Verifizierung der Zeitfolge (vor Buchung/vor Zahlung).

### 14 Menschliche Nachprüfung: Ja
| fehlende Information | Prüfschritt | Prüfhandlung |
|---|---|---|
| Betragsabweichung Freigabenachweis vs. Rechnung | 4/5 Zuordnung, 8 Gesamtbewertung | Klärung mit Genehmiger Person B, auf welchen Betrag sich die Freigabe bezieht |
| Buchungs- und Zahlungsdatum nicht vorhanden | 7 Zeitfolge | Abfrage im Buchhaltungssystem |
| Ursache der Netto-Summenabweichung (5.350 vs. 5.305) | 9 Rechnerische Prüfung | Rückfrage beim Lieferanten / Abgleich Mandatsvereinbarung |

### 15 Audit Exception: Ja
Die Rechnung enthält einen belegten rechnerischen Fehler: Die Summe der Einzelpositionen (5.305,00 €) weicht vom ausgewiesenen Nettobetrag (5.350,00 €) ab, mit Fortpflanzung auf USt und Brutto. Dieselbe Abweichung besteht gegenüber den Soll-Werten der Referenztabelle (+53,55 € brutto). Zusätzlich bezieht sich der Freigabenachweis betragsmäßig auf den Referenz-Soll-Betrag (6.312,95 €), nicht auf das ausgewiesene Rechnungsbrutto (6.366,50 €), was einer eindeutigen Freigabe des tatsächlich geforderten Betrags entgegensteht.

Punkte aus 12 mit „Audit-Exception-relevant: Ja": Nr. 1 und Nr. 2.

GPT 5.5
Antwort 24 (\ac{FN}):

## Rechnung: PA-26-0417 | 24_Rechnung_Patentanmeldung_PA-26-0417.pdf

### 1 Rechnungsdaten
Empfänger: Freudenberg Sealing Technologies GmbH & Co. KG, Abteilung Beschaffung, Höhrenweg 2-4, 69469 Weinheim | Rechnung: Kopfbereich

Rechnungsnr.: PA-26-0417 | Rechnung: Feld „Rechnungsnummer“

Rechnungsdatum: 11.05.2026 | Rechnung: Feld „Rechnungsdatum“

Leistungsdatum: 06.05.2026 | Rechnung: Feld „Leistungsdatum“

Kundennr.: P-2019 | Rechnung: Feld „Kundennummer“

Referenzart/-nr.: Aktenzeichen / DE 10 2026 004 118 | Rechnung: Feld „Aktenzeichen“

Bestellnr.: nicht vorhanden | Rechnung

Lieferant: Patentis Patentanwälte | Rechnung: Briefkopf

Anschriften:
- Lieferant: Erfinderweg 3, 80331 München
- Empfänger: Höhrenweg 2-4, 69469 Weinheim
| Rechnung

St-Nr./USt-IdNr.:
- Lieferant USt-IdNr.: DE 746 517 229
- Lieferant Steuernummer: 19/577/78727
- Empfänger USt-IdNr.: DE 143 566 072
| Rechnung Fußzeile/Kopfbereich

Leistung:
- Ausarbeitung Anmeldeunterlagen (Dichtungssystem, 14 Ansprüche)
- Amtsgebühr DPMA
- Recherchegebühr
- Zeichnungen
| Rechnung Positionen

Menge:
- 1 Pauschale
- 1 Position
- 1 Position
- 9 Blatt

Einzelpreis:
- 4.200,00 €
- 40,00 €
- 300,00 €
- 85,00 €

Positionen:
- 4.200,00 €
- 40,00 €
- 300,00 €
- 765,00 €

Netto: 5.305,00 € | Rechnung

Steuersatz: 19 % | Rechnung

USt: 1.007,95 € | Rechnung

Brutto: 6.312,95 € | Rechnung

Währung: EUR | Rechnung

Zahlungsziel: Zahlbar innerhalb von 21 Tagen ohne Abzug, fällig am 01.06.2026 | Rechnung

IBAN: DE97 6709 0000 9796 9325 09 | Rechnung Fußzeile

Freigabereferenz: Rechnung_PA-26-0417.pdf | Freigabe-Screenshot

### 2 Datensatz
| Datensatz | Kriterium | Werte | Ergebnis |
|---|---|---|---|
| Referenztabelle_Rechnungen.xlsx | Rechnungsnr. exakt | PA-26-0417 | Treffer eindeutig [1] |

### 3 Soll-Ist
| Feld | Ist | Soll | Abweichung abs./% | Status |
|---|---|---|---|---|
| Empfänger | Freudenberg Sealing Technologies GmbH & Co. KG | Freudenberg Sealing Technologies GmbH & Co. KG | 0 / 0 % | Übereinstimmung |
| Rechnungsnr. | PA-26-0417 | PA-26-0417 | 0 / 0 % | Übereinstimmung |
| Rechnungsdatum | 11.05.2026 | 11.05.2026 | 0 / 0 % | Übereinstimmung |
| Leistungsdatum | 06.05.2026 | 06.05.2026 | 0 / 0 % | Übereinstimmung |
| Kundennr. | P-2019 | P-2019 | 0 / 0 % | Übereinstimmung |
| Referenzart | Aktenzeichen | Aktenzeichen | 0 / 0 % | Übereinstimmung |
| Referenznr. | DE 10 2026 004 118 | DE 10 2026 004 118 | 0 / 0 % | Übereinstimmung |
| Lieferant | Patentis Patentanwälte | Patentis Patentanwälte | 0 / 0 % | Übereinstimmung |
| Leistung | Patentanmeldung | Patentanmeldung | 0 / 0 % | Übereinstimmung |
| Netto | 5.305,00 € | 5.305,00 € | 0,00 € / 0 % | Übereinstimmung |
| Steuersatz | 19 % | 19 % | 0 / 0 % | Übereinstimmung |
| USt | 1.007,95 € | 1.007,95 € | 0,00 € / 0 % | Übereinstimmung |
| Brutto | 6.312,95 € | 6.312,95 € | 0,00 € / 0 % | Übereinstimmung |
| Zahlungsziel | 01.06.2026 | 01.06.2026 | 0 / 0 % | Übereinstimmung |
| USt-IdNr. Lieferant | DE 746 517 229 | nicht vorhanden in Referenztabelle | nicht prüfbar | nicht prüfbar |

Sollwerte gemäß Referenztabelle_Rechnungen.xlsx. [1]

### 4 Freigabenachweise (je Nachweis)
| Dokument | Typ | Fundstelle | Absender | Genehmiger | Rolle | Wortlaut | Datum/Uhrzeit | Merkmale (Nr., Lieferant, Betrag) | Lesbarkeit | Evidenz |
|---|---|---|---|---|---|---|---|---|---|---|
| PNG-Screenshot | E-Mail | Betreff und Nachrichtentext | Person B <person.b@freudenberg.example> | Person B | nicht prüfbar | „Rechnung PA-26-0417 (Patentis Patentanwälte) vom 11.05.2026 über 6.312,95 € ist freigegeben.“ | 18.05.2026 15:51 | Rechnungsnr. PA-26-0417; Lieferant Patentis Patentanwälte; Betrag 6.312,95 € | lesbar | hoch |

### 5 Zuordnung
| Nachweis | Merkmale | Werte Rechnung | Werte Nachweis | Ergebnis |
|---|---|---|---|---|
| E-Mail-Screenshot | Rechnungsnr., Lieferant, Betrag | PA-26-0417; Patentis Patentanwälte; 6.312,95 € | PA-26-0417; Patentis Patentanwälte; 6.312,95 € | eindeutig |

### 6 Berechtigung
| Betrag/Betragsart/Währung | Soll (Stufe, Person) | Ist (Genehmiger) | Status | Begründung |
|---|---|---|---|---|
| 6.312,95 € brutto / EUR | Person B, unbegrenzt | Person B | Übereinstimmung | Person B darf jede Rechnung in EUR freigeben. [2] |

### 7 Zeitfolge
| Daten | nach Rechnung | vor Buchung | vor Zahlung (Ja/Nein/Nicht prüfbar) |
|---|---|---|---|
| Rechnungsdatum 11.05.2026; Freigabe 18.05.2026 15:51 | Ja | nicht prüfbar | Nicht prüfbar |

### 8 Gesamtbewertung Freigabe
- Freigabe vorhanden und regelkonform.

### 9 Rechnerische Prüfung
| Werte | berechnet | ausgewiesen | Differenz | Bewertung |
|---|---|---|---|---|
| Positionssumme netto | 5.305,00 € | 5.305,00 € | 0,00 € | korrekt |
| USt 19 % auf 5.305,00 € | 1.007,95 € | 1.007,95 € | 0,00 € | korrekt |
| Brutto | 6.312,95 € | 6.312,95 € | 0,00 € | korrekt |

### 10 Formale Vollständigkeit
| Pflichtangabe | Fundstelle | Status |
|---|---|---|
| Lieferantenanschrift | Briefkopf | vorhanden |
| Leistungsempfänger | Adressfeld | vorhanden |
| Rechnungsdatum | Rechnungsdatenblock | vorhanden |
| Rechnungsnummer | Rechnungsdatenblock | vorhanden |
| Leistungszeitpunkt | Rechnungsdatenblock | vorhanden |
| Leistungsbeschreibung | Positionen | vorhanden |
| Entgelt | Summenbereich | vorhanden |
| Steuersatz | Summenbereich | vorhanden |
| Steuerbetrag | Summenbereich | vorhanden |
| USt-IdNr./Steuernummer | Fußzeile | vorhanden |

### 11 Plausibilität und Betrugsindikatoren
| Kriterium | Werte | plausibel/auffällig/nicht prüfbar |
|---|---|---|
| Leistungsdatum vor Rechnungsdatum | 06.05.2026 vor 11.05.2026 | plausibel |
| Währungskonsistenz | EUR durchgängig | plausibel |
| Steuersatzkonsistenz | 19 % korrekt berechnet | plausibel |
| Empfängerkonsistenz | Rechnung an Freudenberg Sealing Technologies GmbH & Co. KG | plausibel |
| Zahlungsziel | 21 Tage bis 01.06.2026 | plausibel |
| Lieferant in Referenztabelle | vorhanden | plausibel |
| Bestellnummer | nicht vorhanden | auffällig |
| Freigabereferenz | vorhanden (E-Mail-Anhangsbezug) | plausibel |
| Generische Leistung | nein | plausibel |
| Rechnungsdatum Wochenende/Feiertag | 11.05.2026 (Dienstag) | plausibel |
| Freigabedatum Wochenende/Feiertag | 18.05.2026 (Montag) | plausibel |
| IBAN-Inhaber-Abgleich | keine Vergleichsquelle | nicht prüfbar |

### 12 Auffälligkeiten
1. Fehlende Bestellnummer (Bestellnr.: nicht vorhanden) | Risiko: Mittel | Audit-Exception-relevant: Nein | Begründung: Bestellreferenz fehlt auf der Rechnung. | Quelle/Dokument: Rechnung | Prüfdhandlung: Prüfung im Quellsystem auf vorhandene Bestellung oder alternative Beauftragung.

### 13 Empfehlung
- Prüfung im Quellsystem, ob eine Bestellung oder Mandatierung zur Patentanmeldung hinterlegt ist.
- Abgleich der Lieferantenstammdaten (IBAN, USt-IdNr.) mit dem ERP-System.
- Buchungs- und Zahlungsbeleg zur vollständigen Zeitfolgenprüfung einsehen.

### 14 Menschliche Nachprüfung: Ja
- Buchungsdatum nicht vorhanden | Prüfschritt 8 | Buchungsbeleg im Quellsystem prüfen.
- Zahlungsdatum nicht vorhanden | Prüfschritt 8 | Zahlungsbeleg im Quellsystem prüfen.
- USt-IdNr. des Lieferanten nicht in Referenztabelle enthalten | Prüfschritt 3 | Stammdatenabgleich im Quellsystem durchführen.

### 15 Audit Exception: Nein
Es wurde ein eindeutiger Referenzdatensatz gefunden, die Soll-Ist-Werte stimmen überein, die Rechnung ist rechnerisch korrekt und eine eindeutige sowie berechtigte Freigabe durch Person B liegt vor. [1]​[2] Die festgestellte fehlende Bestellnummer ist gemäß Vorgaben nicht Audit-Exception-relevant.

Opus 4.7
Antwort 98 (\ac{TP})

Die Referenztabelle enthält keine USt-IdNr. des Lieferanten als separates Feld. Ich habe alle notwendigen Informationen. Nun erstelle ich das Prüfprotokoll.

## Rechnung: FW-2026-0771 | 98_Rechnung_Sommerfest_FW-2026-0771.pdf

### 1 Rechnungsdaten
- **Empfänger:** Freudenberg & Co. Kommanditgesellschaft, Zentraler Einkauf, Höhnerweg 2-4, 69469 Weinheim
- **Rechnungsnummer:** FW-2026-0771
- **Rechnungsdatum:** 04.06.2026
- **Leistungsdatum:** 30.05.2026
- **Kundennummer:** FE-3311
- **Referenzart/-nr.:** Veranstaltung / Sommerfest 2026
- **Bestellnummer:** nicht vorhanden
- **Lieferant:** Festwerk Eventagentur GmbH, Festplatzweg 7, 68159 Mannheim
- **USt-IdNr. Empfänger:** DE 143 566 072
- **USt-IdNr. Lieferant:** DE 924 942 840 (Fußzeile)
- **St-Nr. Lieferant:** 75/368/97848 (Fußzeile)
- **Leistung:** Abrechnung Betriebssommerfest vom 30.05.2026 (Catering/Bewirtung, Bühne/Ton/Licht, Zelte/Mobiliar, Rahmenprogramm)
- **Positionen:**
  - Pos. 1: 1.800 Gäste × 38,00 € = 68.400,00 €
  - Pos. 2: 1 Pauschale × 14.800,00 € = 14.800,00 €
  - Pos. 3: 1 Pauschale × 9.600,00 € = 9.600,00 €
  - Pos. 4: 1 Pauschale × 6.800,00 € = 6.800,00 €
- **Netto:** 99.600,00 €
- **Steuersatz:** 19 %
- **USt:** 18.924,00 €
- **Brutto:** 118.524,00 €
- **Währung:** EUR
- **Zahlungsziel:** 21 Tage, fällig am 25.06.2026
- **IBAN:** DE24 6001 0070 7569 4185 80 (Postbank Stuttgart, BIC PBNKDEFFXXX)
- **Freigabereferenz:** nicht vorhanden

### 2 Datensatz
| Datensatz | Kriterium | Werte | Ergebnis |
|---|---|---|---|
| 98_Rechnung_Sommerfest_FW-2026-0771.pdf[1] | Rechnungsnr. exakt | FW-2026-0771 | Treffer |
| | Lieferant | Festwerk Eventagentur GmbH | Übereinstimmung |
| | Brutto | 118.524,00 € | Übereinstimmung |
| | Rechnungsdatum | 04.06.2026 | Übereinstimmung |

Eindeutiger Referenzdatensatz gefunden.

### 3 Soll-Ist
| Feld | Ist | Soll | Abweichung abs./% | Status |
|---|---|---|---|---|
| Empfänger | Freudenberg & Co. Kommanditgesellschaft | Freudenberg & Co. Kommanditgesellschaft[1] | – | Übereinstimmung |
| Rechnungsnr. | FW-2026-0771 | FW-2026-0771[1] | – | Übereinstimmung |
| Rechnungsdatum | 04.06.2026 | 04.06.2026[1] | – | Übereinstimmung |
| Leistungsdatum | 30.05.2026 | 30.05.2026 (Leistungsdatum)[1] | – | Übereinstimmung |
| Kundennummer | FE-3311 | FE-3311[1] | – | Übereinstimmung |
| Referenzart | Veranstaltung | Veranstaltung[1] | – | Übereinstimmung |
| Referenznr. | Sommerfest 2026 | Sommerfest 2026[1] | – | Übereinstimmung |
| Lieferant | Festwerk Eventagentur GmbH | Festwerk Eventagentur GmbH[1] | – | Übereinstimmung |
| Leistung | Sommerfest/Betriebsfeier | Sommerfest Betriebsfeier[1] | – | Übereinstimmung |
| Netto | 99.600,00 € | 99,600.00 €[1] | 0,00 € / 0,00 % | Übereinstimmung |
| Steuersatz | 19 % | 19%[1] | – | Übereinstimmung |
| USt | 18.924,00 € | 18,924.00 €[1] | 0,00 € / 0,00 % | Übereinstimmung |
| Brutto | 118.524,00 € | 118,524.00 €[1] | 0,00 € / 0,00 % | Übereinstimmung |
| Zahlungsziel/fällig | 25.06.2026 | 25.06.2026[1] | – | Übereinstimmung |
| USt-IdNr. Lieferant | DE 924 942 840 | nicht vorhanden in Referenztabelle | – | nicht prüfbar |

### 4 Freigabenachweise
| Dokument | Typ | Fundstelle | Absender | Genehmiger | Rolle | Wortlaut | Datum/Uhrzeit | Merkmale | Lesbarkeit | Evidenz |
|---|---|---|---|---|---|---|---|---|---|---|
| PNG-Screenshot E-Mail | E-Mail-Antwort | Betreff „AW: Rechnung FW-2026-0771 zur Freigabe" | Person E <person.e@freudenberg.example> | Person E | nicht angegeben | „die Rechnung FW-2026-0771 der Festwerk Eventagentur GmbH über 118.524,00 € ist freigegeben. Die Leistung wurde wie beauftragt erbracht." | Di 09.06.2026 11:38 | Rechnungsnr. FW-2026-0771; Lieferant Festwerk Eventagentur GmbH; Betrag 118.524,00 € | lesbar | hoch (Genehmiger, Inhalt, Datum, Rechnungsbezug) |

### 5 Zuordnung
| Nachweis | Merkmale | Werte Rechnung | Werte Nachweis | Ergebnis |
|---|---|---|---|---|
| E-Mail Person E | Rechnungsnr. exakt | FW-2026-0771 | FW-2026-0771 | eindeutig |
| | Lieferant | Festwerk Eventagentur GmbH | Festwerk Eventagentur GmbH | eindeutig |
| | Brutto | 118.524,00 € | 118.524,00 € | eindeutig |

Status: **eindeutig zugeordnet**.

### 6 Berechtigung
| Betrag/Betragsart/Währung | Soll (Stufe, Person) | Ist (Genehmiger) | Status | Begründung |
|---|---|---|---|---|
| 118.524,00 € Brutto / EUR | Person E: keine Berechtigung, Darf keine Rechnung freigeben[2] | Person E | nicht berechtigt | Person E hat laut Freigabematrix generell keine Freigabeberechtigung. |

### 7 Zeitfolge
| Daten | nach Rechnung | vor Buchung | vor Zahlung |
|---|---|---|---|
| Rechnungsdatum 04.06.2026; Freigabe 09.06.2026; Fälligkeit 25.06.2026 | Ja (09.06. nach 04.06.) | Nicht prüfbar (Buchungsdatum liegt nicht vor) | Nicht prüfbar (Zahlungsdatum liegt nicht vor; Freigabe liegt vor Fälligkeit) |

### 8 Gesamtbewertung Freigabe
- **Freigabe vorhanden, Genehmiger nicht berechtigt.**

### 9 Rechnerische Prüfung
| Werte | berechnet | ausgewiesen | Differenz | Bewertung |
|---|---|---|---|---|
| Pos. 1: 1.800 × 38,00 € | 68.400,00 € | 68.400,00 € | 0,00 € | korrekt |
| Pos. 2 | 14.800,00 € | 14.800,00 € | 0,00 € | korrekt |
| Pos. 3 | 9.600,00 € | 9.600,00 € | 0,00 € | korrekt |
| Pos. 4 | 6.800,00 € | 6.800,00 € | 0,00 € | korrekt |
| Netto (Σ Positionen) | 99.600,00 € | 99.600,00 € | 0,00 € | korrekt |
| USt (99.600 × 19 %) | 18.924,00 € | 18.924,00 € | 0,00 € | korrekt |
| Brutto (Netto + USt) | 118.524,00 € | 118.524,00 € | 0,00 € | korrekt |

### 10 Formale Vollständigkeit
| Pflichtangabe | Fundstelle | Status |
|---|---|---|
| Anschrift Lieferant | Kopf/Fuß | vorhanden |
| Anschrift Empfänger | Rechnungsempfänger | vorhanden |
| St-Nr. Lieferant | Fußzeile | vorhanden |
| USt-IdNr. Lieferant | Fußzeile (DE 924 942 840) | vorhanden |
| USt-IdNr. Empfänger | Empfängerblock (DE 143 566 072) | vorhanden |
| Rechnungsdatum | Rechnungsdaten | vorhanden |
| Rechnungsnummer | Rechnungsdaten / Überschrift | vorhanden |
| Leistungsbeschreibung | Positionstabelle | vorhanden |
| Leistungszeitpunkt | Rechnungsdaten (30.05.2026) | vorhanden |
| Entgelt (Netto) | Summenblock | vorhanden |
| Steuersatz | Summenblock (19 %) | vorhanden |
| Steuerbetrag | Summenblock | vorhanden |

### 11 Plausibilität und Betrugsindikatoren
| Kriterium | Werte | Bewertung |
|---|---|---|
| Leistungs- nicht nach Rechnungsdatum | Leistung 30.05.2026, Rechnung 04.06.2026 | plausibel |
| Währung konsistent | EUR durchgehend | plausibel |
| Steuersatz konsistent | 19 % durchgehend | plausibel |
| Empfänger konsistent | Freudenberg & Co. KG, Weinheim | plausibel |
| IBAN-Inhaber-Vergleich | keine Vergleichsdaten | nicht prüfbar |
| Zahlungsziel | 21 Tage, fällig 25.06.2026 | plausibel |
| Fehlende Bestellnummer | – | auffällig (Red Flag 11d) |
| Fehlende Freigabereferenz auf Rechnung | – | auffällig (Red Flag 11d) |
| Freigabe über externe Adresse | person.e@freudenberg.example (interne Domain) | plausibel |
| Rechnungs-/Freigabedatum Wochenende/Feiertag | 04.06.2026 Do; 09.06.2026 Di | plausibel |
| Marktüblichkeit der Beträge | – | nicht prüfbar |
| Lieferant in Referenztabelle | Treffer FW-2026-0771 | plausibel |

### 12 Auffälligkeiten
1. **Freigabe durch nicht berechtigte Person (Person E, 118.524,00 € Brutto)** | Risiko: Hoch (Freigabe unberechtigt, materieller Betrag) | **Audit-Exception-relevant: Ja** | Begründung: Person E: keine Berechtigung, Darf keine Rechnung freigeben[2]; Rechnung wurde dennoch per E-Mail vom 09.06.2026 freigegeben. | Quelle: Freigabematrix.xlsx; PNG-Screenshot „AW: Rechnung FW-2026-0771 zur Freigabe". | Prüfhandlung: Klärung im Fachbereich, nachträgliche Freigabe durch berechtigte Person (A oder B) einholen; Vier-Augen-Kontrolle im Freigabeprozess verifizieren.
2. **Fehlende Bestellnummer auf der Rechnung** | Risiko: Mittel (fehlende Bestellreferenz) | Audit-Exception-relevant: Nein | Begründung: Auf der Rechnung ist keine Bestellnummer ausgewiesen; Referenzart laut Referenztabelle ist „Veranstaltung". | Quelle: 98_Rechnung_Sommerfest_FW-2026-0771.pdf. | Prüfhandlung: Prüfen, ob eine Bestellung/Auftragsvergabe im Quellsystem existiert.

### 13 Empfehlung
- Nachträgliche Freigabe durch eine gemäß Freigabematrix berechtigte Person (Person A oder Person B) einholen; vorherige Freigabe durch Person E ist unwirksam.
- Im Quellsystem prüfen, ob die Rechnung bereits gebucht oder zur Zahlung angewiesen wurde; ggf. Zahlungslauf stoppen.
- Vorliegen einer Bestellung / eines Auftrags zur Veranstaltung „Sommerfest 2026" (FE-3311) im Quellsystem prüfen.
- Abgleich IBAN DE24 6001 0070 7569 4185 80 mit Lieferantenstammdaten im Kreditorensystem.
- Berechtigungsprozess überprüfen: Wie konnte eine Anfrage zur Freigabe an eine nicht berechtigte Person (Person E) gelangen bzw. akzeptiert werden.

### 14 Menschliche Nachprüfung: Ja
| fehlende Information | Prüfschritt | Prüfhandlung |
|---|---|---|
| Nachweisbare Freigabe durch berechtigte Person | 4/6 Freigabe & Berechtigung | Berechtigte Freigabe (Person A/B) einholen und dokumentieren |
| Buchungs- und Zahlungsstatus | 7 Zeitfolge | Status im Quellsystem prüfen, ggf. Zahlung stoppen |
| IBAN-Abgleich mit Stammdaten | 11c Plausibilität | Abgleich im Kreditorensystem |

### 15 Audit Exception: Ja
Die Rechnung wurde ausschließlich durch Person E freigegeben, die laut Freigabematrix generell nicht freigabeberechtigt ist. Damit liegt eine fehlerhafte Freigabe vor; der Betrag von 118.524,00 € Brutto übersteigt zudem die Grenzen von Person C (25.000,00 €) und Person D (1.000,00 €), sodass nur Person A oder B befugt gewesen wären.

Audit-Exception-relevante Punkte aus Abschnitt 12:
1. Freigabe durch nicht berechtigte Person (Person E).


GPT 5.5
Antwort 98 (\ac{TP}):

## Rechnung: FW-2026-0771 | 98_Rechnung_Sommerfest_FW-2026-0771.pdf

### 1 Rechnungsdaten
Empfänger, Rechnungsnr., -datum, Leistungsdatum, Kundennr., Referenzart/-nr., Bestellnr., Lieferant, Anschriften, St-Nr./USt-IdNr., Leistung, Menge, Einzelpreis, Positionen, Netto, Steuersatz, USt, Brutto, Währung, Zahlungsziel, IBAN, Freigabereferenz

Empfänger: Freudenberg & Co. Kommanditgesellschaft, Zentraler Einkauf, Höhnerweg 2-4, 69469 Weinheim

Rechnungsnr.: FW-2026-0771

Rechnungsdatum: 04.06.2026

Leistungsdatum: 30.05.2026

Kundennr.: FE-3311

Referenzart/-nr.: Veranstaltung / Sommerfest 2026

Bestellnr.: nicht vorhanden

Lieferant: Festwerk Eventagentur GmbH

Anschriften:
- Lieferant: Festplatzweg 7, 68159 Mannheim
- Empfänger: Höhnerweg 2-4, 69469 Weinheim

St-Nr./USt-IdNr.:
- Rechnungskopf: USt-IdNr. DE 143 566 072
- Fußzeile: USt-IdNr. DE 924 942 840
- St.-Nr. nicht vorhanden

Leistung:
- Catering, Bewirtung vor Ort, Buffet und Getränke, Servicepersonal
- Bühne, Ton- und Lichttechnik
- Zelt und Mobiliar
- Rahmenprogramm und Künstler

Menge:
- 1.800 Gäste
- 1 Pauschale
- 1 Pauschale
- 1 Pauschale

Einzelpreis:
- 38,00 €
- 14.800,00 €
- 9.600,00 €
- 6.800,00 €

Positionen:
- 68.400,00 €
- 14.800,00 €
- 9.600,00 €
- 6.800,00 €

Netto: 99.600,00 €

Steuersatz: 19 %

USt: 18.924,00 €

Brutto: 118.524,00 €

Währung: EUR

Zahlungsziel: innerhalb von 21 Tagen, fällig am 25.06.2026

IBAN: DE24 6001 0007 7569 4185 80

Freigabereferenz: nicht vorhanden

### 2 Datensatz

| Datensatz | Kriterium | Werte | Ergebnis |
|---|---|---|---|
| Referenztabelle_Rechnungen.xlsx | Rechnungsnr. exakt | FW-2026-0771 | Eindeutiger Treffer [1] |

### 3 Soll-Ist

| Feld | Ist | Soll | Abweichung abs./% | Status |
|---|---|---|---|---|
| Empfänger | Freudenberg & Co. Kommanditgesellschaft | Freudenberg & Co. Kommanditgesellschaft | 0 / 0 % | Übereinstimmung |
| Rechnungsnr. | FW-2026-0771 | FW-2026-0771 | 0 / 0 % | Übereinstimmung |
| Rechnungsdatum | 04.06.2026 | 04.06.2026 | 0 / 0 % | Übereinstimmung |
| Leistungsdatum | 30.05.2026 | 30.05.2026 (Leistungsdatum) | 0 / 0 % | Übereinstimmung |
| Kundennr. | FE-3311 | FE-3311 | 0 / 0 % | Übereinstimmung |
| Referenzart | Veranstaltung | Veranstaltung | 0 / 0 % | Übereinstimmung |
| Referenznr. | Sommerfest 2026 | Sommerfest 2026 | 0 / 0 % | Übereinstimmung |
| Lieferant | Festwerk Eventagentur GmbH | Festwerk Eventagentur GmbH | 0 / 0 % | Übereinstimmung |
| Leistung | Abrechnung des Betriebssommerfestes vom 30.05.2026 | Sommerfest Betriebsfeier | nicht berechenbar | Übereinstimmung |
| Netto | 99.600,00 € | 99.600,00 € | 0,00 € / 0 % | Übereinstimmung |
| Steuersatz | 19 % | 19 % | 0 / 0 % | Übereinstimmung |
| USt | 18.924,00 € | 18.924,00 € | 0,00 € / 0 % | Übereinstimmung |
| Brutto | 118.524,00 € | 118.524,00 € | 0,00 € / 0 % | Übereinstimmung |
| Zahlungsziel/Fälligkeit | 25.06.2026 | 25.06.2026 | 0 / 0 % | Übereinstimmung |
| USt-IdNr. Lieferant | DE 143 566 072 / DE 924 942 840 | nicht vorhanden | nicht prüfbar | nicht prüfbar |

### 4 Freigabenachweise (je Nachweis)

| Dokument | Typ | Fundstelle | Absender | Genehmiger | Rolle | Wortlaut | Datum/Uhrzeit | Merkmale (Nr., Lieferant, Betrag) | Lesbarkeit | Evidenz |
|---|---|---|---|---|---|---|---|---|---|---|
| PNG-Screenshot E-Mail | E-Mail | Antwortnachricht | Person E <person.e@freudenberg.example> | Person E | nicht vorhanden | „die Rechnung FW-2026-0771 der Festwerk Eventagentur GmbH über 118.524,00 € ist freigegeben. Die Leistung wurde wie beauftragt erbracht.“ | 09.06.2026 11:38 | Rechnungsnr. FW-2026-0771, Lieferant Festwerk Eventagentur GmbH, Betrag 118.524,00 € | lesbar | hoch |

### 5 Zuordnung

| Nachweis | Merkmale | Werte Rechnung | Werte Nachweis | Ergebnis |
|---|---|---|---|---|
| E-Mail Freigabe | Rechnungsnr., Lieferant, Betrag | FW-2026-0771; Festwerk Eventagentur GmbH; 118.524,00 € | FW-2026-0771; Festwerk Eventagentur GmbH; 118.524,00 € | eindeutig |

### 6 Berechtigung

| Betrag/Betragsart/Währung | Soll (Stufe, Person) | Ist (Genehmiger) | Status | Begründung |
|---|---|---|---|---|
| 118.524,00 € brutto / EUR | Person A oder Person B (unbegrenzt) [2] | Person E | Abweichung | Person E besitzt keine Freigabeberechtigung. [2] |

### 7 Zeitfolge

| Daten | nach Rechnung | vor Buchung | vor Zahlung (Ja/Nein/Nicht prüfbar) |
|---|---|---|---|
| Freigabe 09.06.2026 11:38, Rechnung 04.06.2026 | Ja | nicht prüfbar | Nicht prüfbar |

### 8 Gesamtbewertung Freigabe

- Freigabe vorhanden, Genehmiger nicht berechtigt.

### 9 Rechnerische Prüfung

| Werte | berechnet | ausgewiesen | Differenz | Bewertung |
|---|---|---|---|---|
| 1.800 × 38,00 € | 68.400,00 € | 68.400,00 € | 0,00 € | korrekt |
| Summe Positionen | 99.600,00 € | 99.600,00 € | 0,00 € | korrekt |
| 99.600,00 € × 19 % | 18.924,00 € | 18.924,00 € | 0,00 € | korrekt |
| Netto + USt | 118.524,00 € | 118.524,00 € | 0,00 € | korrekt |

### 10 Formale Vollständigkeit

| Pflichtangabe | Fundstelle | Status |
|---|---|---|
| Name/Anschrift Lieferant | Rechnungskopf/Fußzeile | vorhanden |
| Name/Anschrift Empfänger | Rechnungskopf | vorhanden |
| Rechnungsnummer | Rechnungsdaten | vorhanden |
| Rechnungsdatum | Rechnungsdaten | vorhanden |
| Leistungsbeschreibung | Positionsliste | vorhanden |
| Leistungszeitpunkt | Rechnungsdaten | vorhanden |
| Entgelt | Summenbereich | vorhanden |
| Steuersatz | Summenbereich | vorhanden |
| Steuerbetrag | Summenbereich | vorhanden |
| USt-IdNr./St.-Nr. | Rechnungskopf/Fußzeile | vorhanden |

### 11 Plausibilität und Betrugsindikatoren

| Kriterium | Werte | plausibel/auffällig/nicht prüfbar |
|---|---|---|
| Leistungsdatum vor Rechnungsdatum | 30.05.2026 vor 04.06.2026 | plausibel |
| Währung konsistent | EUR | plausibel |
| Zahlungsziel konsistent | Rechnung 04.06.2026, Fälligkeit 25.06.2026 | plausibel |
| Lieferant in Referenztabelle | vorhanden [1] | plausibel |
| USt-IdNr. innerhalb Rechnung | DE 143 566 072 und DE 924 942 840 | auffällig |
| Fehlende Bestellnummer | keine Bestellnummer angegeben | auffällig |
| Generische Leistung | Sommerfest/Betriebsfeier | auffällig |

### 12 Auffälligkeiten

1. Freigabe durch Person E bei Rechnungsbetrag 118.524,00 € brutto | Hoch | Audit-Exception-relevant: Ja | Genehmiger laut Freigabematrix nicht freigabeberechtigt. [2] | Freigabematrix.xlsx, E-Mail-Screenshot | Berechtigungsprüfung gegen Freigabematrix

2. Zwei unterschiedliche USt-IdNr. auf derselben Rechnung (DE 143 566 072 / DE 924 942 840) | Niedrig | Audit-Exception-relevant: Ja | Belegte Abweichung innerhalb des Rechnungsdokuments. | Rechnung Kopf- und Fußzeile | Dokumentenprüfung

3. Fehlende Bestellnummer | Mittel | Audit-Exception-relevant: Nein | Bestellreferenz nicht angegeben. | Rechnung | Dokumentenprüfung

### 13 Empfehlung

- Berechtigte Freigabe gemäß Freigabematrix im Quellsystem prüfen.
- Ursprung der unterschiedlichen USt-IdNr. beim Lieferanten klären.
- Buchungs- und Zahlungsdaten im Quellsystem prüfen.
- Bestellbezug bzw. Beschaffungsvorgang im Quellsystem nachvollziehen.

### 14 Menschliche Nachprüfung: Ja

- Buchungsdatum fehlt | Prüfschritt 8 | Buchungsbeleg im Quellsystem prüfen.
- Zahlungsdatum fehlt | Prüfschritt 8 | Zahlungsbeleg im Quellsystem prüfen.
- USt-IdNr. des Referenzdatensatzes fehlt | Prüfschritt 3 | Stammdaten/Referenztabelle prüfen.

### 15 Audit Exception: Ja

Ausschlaggebend sind die belegte Freigabe durch Person E trotz fehlender Freigabeberechtigung gemäß Freigabematrix sowie die belegte Abweichung durch zwei unterschiedliche USt-IdNr. auf derselben Rechnung. Die rechnerische Prüfung und der Soll-Ist-Abgleich der übrigen prüfbaren Felder ergaben keine Abweichungen.

